% interface=en output=pdftex

\loadmapfile[texnansi-urw-zapfchan]
\loadmapfile[ec-urw-zapfchan]
\loadmapfile[8r-urw-zapfchan]

\environment mstyle 

\starttext 

\TitlePage{513}{0}{\TeX FONT\\explained}{}{1}

\subject {Introduction}

It is probably a known fact that \TEX\ can handle fonts of any kind, as well
as satisfy the needs of those who do not live in english speaking countries.
Nevertheless, the majority of documents typeset by \TEX, use the more or less
standard Computer Modern Typefaces in an encoding not that well suited for
non english usage. 

The (subtle) relations between font encodings, hyphenation patterns, input
encodings, operating systems, font metrics, and alike makes font handling the
most complicated (and fuzzy) part of macro packages. 

Although \TEX\ distributions come with many fonts, in often obscure names,
made even more obscure by encoding specific instances and|/|or exceptions.
Browsing the font sub||tree of the \type {texmf} directory can drive you
crazy. 

As a result, installing a font, or at least making sure that the fonts
installed already can be used, is not trivial. It is not enough to make your
macro package aware of their existence, you also should inform the post
processor (which in the case of \PDFTEX\ is build in) of (re|)|encoding
issues of shape manipulations. And of course each back||end has its own
demands. 

In \CONTEXT, font support is implemented in layers. One can use an individual
font, define a relationship between them, and collect such relationships in
typeface collections. These issues are covered in the font manual that comes
with \CONTEXT. 

In this manual we will describe \TEXFONT, a \PERL\ script that can help you
to manage the files that are needed for \TEX\ and the back||ends. We hope
that this tool is of help, although we cannot guarantee that life will be
easy from now on. The script only covers \quote {normal} fonts. For setting
up math fonts, vendors often provide the necessary files and \CONTEXT\ is
aware of these files. Special fonts, like expert fonts, assume a more in
depth knowledge of font handling. We may deal with them in the future. 

The more demanding user can of course fall back on more complicated tools
like \type {fontinst}. Although I never used this tool, my impression is that
it does a good job on manipulating fonts, but most users will not need this
kind of fine tuning. 

{\em Currently this manual is directed towards \CONTEXT\ users who also use 
\PDFTEX. Future versions may also cover \DVI\ (to \PDF) drivers.} 

\subject {Organizing fonts}

In the \TEX\ community, much comes for free. This is definitely not true for
fonts, although a decent free collection is available to get you started.
There are thousands of interesting fonts out there, and occasionally you may
want to buy one. 

So, what to do when this floppy or \CDROM\ arrives. A quick look at the
contents will show you that there are at least \type {pfb} and \type {afm}
files on it. If not, then you have a problem. But, given that these files are
there, where should they go? 

Regular \TEX\ distributions are organized conforming the \TEX\ Directory
Structure. There you will find the commonly used and stable components under
the main tree. In this tree you will find a font sub||tree: 

\starttyping
tex/texmf/fonts
\stoptyping

Any deviation from the default components, being your own extensions or
updates, should go into a local tree: 

\starttyping
tex/texmf-local/fonts
\stoptyping

The main tree is preferably just a copy of for instance the one that comes on
\TEX\ Live. The \type {tex/texmf-local} path is where you normally will unzip
a \CONTEXT\ update. This is also a natural place to put your fonts, given
that you have access to this tree. That way, you can easily replace the main
tree without spoiling your local settings. 

When you buy a font for your own usage only, the local tree is a good place
for them. But when you want to share them you need to be aware of licensing
issues. Licences may permit installation on 5~machines, printing on
1~printer, placement on a server with restricted access for 3 simultaneous
users, and alike. 

So, this is why it may make sense to introduce another tree, which we can
maintain separately. In this tree we put those files that are needed to use
the fresh fonts. So we may have at least the following paths: 

\starttyping
tex/texmf-fonts/fonts
tex/texmf-fonts/pdftex/config
\stoptyping

Below the \type {/fonts} subpath the font specific files are organized 
according to vendor and collection. 

\starttyping 
texmf / fonts  / tfm    / vendor   / name / *.tfm 
               / afm    / vendor   / name / *.afm
               / pfm    / vendor   / name / *.pfm
               / vf     / vendor   / name / *.vf 
               / type1  / vendor   / name / *.pfb 
\stoptyping 

The \PDFTEX\ directory is organized as follows 

\starttyping 
texmf / pdftex / config /                   *.cfg 
               / config /                   *.map
               / config / encoding /        *.enc
\stoptyping 

Of course you need to set up a couple of environment variables in order to
let your \TEX\ system understand this all. Here we assume that you use \WEBC. 

\starttyping
TEXMFMAIN=/tex/texmf
TEXMFLOCAL=/tex/texmf-local
TEXMFFONTS=/tex/texmf-fonts
\stoptyping

The order of searching these trees is determined by the following variable: 

\starttyping
TEXMF={$TEXMFFONTS,$TEXMFLOCAL,!!$TEXMFMAIN}
\stoptyping

Specific settings can be taken care of in the \type {texmf.cnf} file. The
previous definition could have gone there. 

\starttyping
TEXMFCNF=/tex/texmf-local/web2c
\stoptyping

Since \TEXFONT\ looks for both the local and dedicated font tree, you're free
in your choice, although from the perspective of managing fonts it may be
handy to use the dedicated tree. 

Back to the question \quotation {Where should my font files go?} we can now
answer: \quotation {In the font tree.}. Say that you've bought Officina from
ITC, you can put the files in: 

\starttyping
/tex/texmf-fonts/fonts/source/itc/officina
\stoptyping

The files that \TEX\ needs for using these fonts, will later be moved to
and|/|or put in: 

\starttyping
/tex/texmf-fonts/fonts/afm/itc/officina
/tex/texmf-fonts/fonts/tfm/itc/officina
/tex/texmf-fonts/fonts/vf/itc/officina
/tex/texmf-fonts/fonts/type1/itc/officina
\stoptyping

Here, \type {type1} does not denote a suffix, but a more verbose naming of
\type {pfb} files. The \type {tfm} files contain the metric information that
\TEX\ needs during the typesetting process, while the virtual font files
\type {vf} are needed in the postprocessing stage to sort out how glyphs are
composed (if they are composed at all). 

\subject {Generating metrics} 

Before you can set up \CONTEXT\ to use fonts, you need to generate the metric
files needed. 

\starttyping
texfont --ve=itc --co=officina --ma --in 
\stoptyping

This command means as much as: we're going to handle the Officina collection
from vendor ITC, make all directories needed, and install the files there.
Installation here comes down to copying the \type {afm} and \type {pfb}
files, as well as generating the \type {tfm} and \type {vf} files. 

We started in the installation path with files like: 

\starttyping
ovb_____.afm
ovb_____.pfb
\stoptyping

and end up with files spread over the font tree named: 

\starttyping
/tex/fonts/afm/itc/officina/ovb_____.afm
/tex/fonts/type1/itc/officina/ovb_____.pfb
/tex/fonts/tfm/itc/officina/raw-ovb.tfm
/tex/fonts/tfm/itc/officina/texnansi-ovb.tfm
/tex/fonts/vf/itc/officina/texnansi-raw-ovb.tfm
\stoptyping

From this you can deduce that we clean up names, as well as handle an
encoding vector. By default we use the \type {texnansi} vector, but you can
specify another one if needed: 

\starttyping
texfont --en=ec --ve=itc --co=officina --ma --in 
\stoptyping

\TEXFONT\ also generates a \TEX\ file that demonstrates the usage of these
fonts. Normally you will define typescripts to take care of this, or use the
predefined ones that come with \CONTEXT. 

\starttyping
\definefontsynonym[OfficinaSerif-Bold][ec-ovb][encoding=ec]
\stoptyping

We also need to instruct \PDFTEX\ on how to embed the font. For this purpose,
a map file is generated: 

\starttyping
\loadmapfile[ec-itc-officina.map]
\stoptyping

In this map file, you will find lines like: 

\starttyping
ec-raw-ovb OfficinaSerif-Bold 4 < ovb_____.pfb ec.enc
\stoptyping

In practice it means that \TEX\ will base its typesetting decisions on the
metric file, but for inclusion will fall back on the associated (to be
reencoded) file. This file has references to the raw file and these
references are resolved in the map file. 

The default encoding is \type {texnansi}. For western languages, \type {8r}
and \type {ec} are also a good choice. But, in any case, make sure that all
the characters that you need are there by processing the generated \type
{tex} file. 

\starttyping
texexec ec-itc-officina
\stoptyping

If this file processes all right, then at least you know that you have the
correct files on your system. When processing this file, you can enable compact 
mode. 

\starttyping
texexec ec-itc-officina --mode=compact 
\stoptyping

\subject {Example}

Normally you need more then one run of \TEXFONT\ to create the files needed.
The next sequence removes old instances and installs the new fonts. We assume
that you have put the \type {afm} and \type {pfb} files on the sourcepath
relative to the current location. 

\starttyping
texfont --ve=itc --co=officina --re 
texfont --ve=itc --co=officina --so=itc/officina --ma --in

texfont --ve=itc --co=officina --so=itc/officina --ca=* ovbk_*
texfont --ve=itc --co=officina --so=itc/officina --sl=* ovbk_*
texfont --ve=itc --co=officina --so=itc/officina --sl=* ovb_*

texfont --ve=itc --co=officina --so=itc/officina --ca=* owbk_*
texfont --ve=itc --co=officina --so=itc/officina --sl=* owbk_*
texfont --ve=itc --co=officina --so=itc/officina --sl=* owb_*
\stoptyping

For setting up \CONTEXT\ to actually use these fonts, you need to take a look
at the \type {tex} file that is generated and|/|or take a look at the font
manual. Of course you can ask around to see if someone already has the
typescripts made. Without comment we show how such a definition looks. We
assume that you put these definitions in a file called \type {typeface.tex}. 

\start \setuptyping[bodyfont=small,blank=medium]

\starttyping
\starttypescript [map] [8r]

\loadmapfile [8r-itc-officina]

\stoptypescript

\starttypescript [serif] [officina] [name]

\definefontsynonym[Serif]           [OfficinaSerif-Book]
\definefontsynonym[SerifItalic]     [OfficinaSerif-BookItalic]
\definefontsynonym[SerifSlanted]    [OfficinaSerif-BookSlanted]
\definefontsynonym[SerifBold]       [OfficinaSerif-Bold]
\definefontsynonym[SerifBoldItalic] [OfficinaSerif-BoldItalic]
\definefontsynonym[SerifBoldSlanted][OfficinaSerif-BoldSlanted]
\definefontsynonym[SerifCaps]       [OfficinaSerif-Caps]

\stoptypescript

\starttypescript [sans] [officina] [name]

\definefontsynonym[Sans]           [OfficinaSans-Book]
\definefontsynonym[SansItalic]     [OfficinaSans-BookItalic]
\definefontsynonym[SansSlanted]    [OfficinaSans-BookSlanted]
\definefontsynonym[SansBold]       [OfficinaSans-Bold]
\definefontsynonym[SansBoldItalic] [OfficinaSans-BoldItalic]
\definefontsynonym[SansBoldSlanted][OfficinaSans-BoldSlanted]
\definefontsynonym[SansCaps]       [OfficinaSans-Caps]

\stoptypescript

\starttypescript [serif] [officina] [8r]

\definefontsynonym[OfficinaSerif-Book]       [8r-ovbk][encoding=8r]
\definefontsynonym[OfficinaSerif-BookItalic] [8r-ovwi][encoding=8r]
\definefontsynonym[OfficinaSerif-Bold]       [8r-ovb] [encoding=8r]
\definefontsynonym[OfficinaSerif-BoldItalic] [8r-ovbi][encoding=8r]

\definefontsynonym[OfficinaSerif-BookSlanted][8r-ovbk-slanted-167]    [encoding=8r]
\definefontsynonym[OfficinaSerif-BoldSlanted][8r-ovb-slanted-167]     [encoding=8r]
\definefontsynonym[OfficinaSerif-Caps]       [8r-ovbk-capitalized-800][encoding=8r]

\stoptypescript

\starttypescript [sans] [officina] [8r]

\definefontsynonym[OfficinaSans-Book]       [8r-owbk]                [encoding=8r]
\definefontsynonym[OfficinaSans-BookItalic] [8r-owwi]                [encoding=8r]
\definefontsynonym[OfficinaSans-Bold]       [8r-owb]                 [encoding=8r]
\definefontsynonym[OfficinaSans-BoldItalic] [8r-owbi]                [encoding=8r]
\definefontsynonym[OfficinaSans-BookSlanted][8r-owbk-slanted-167]    [encoding=8r]
\definefontsynonym[OfficinaSans-BoldSlanted][8r-owb-slanted-167]     [encoding=8r]
\definefontsynonym[OfficinaSans-Caps]       [8r-owbk-capitalized-800][encoding=8r]

\stoptypescript
\stoptyping

\stop

In your document style, you can now stick to simple definitions like: 

\starttyping 
\usetypescriptfile[typeface]

\definetypeface[officina][rm][serif][officina][default][encoding=8r]
\definetypeface[officina][ss][sans] [officina][default][encoding=8r]

\setupbodyfont[officina,rm,10pt]
\stoptyping 

Of course you can mix these fonts with other ones, in which case you may want
to apply relative scaling first. The fonts manual also explains how you can
set up this font to use protruding characters. Since we use symbolic names,
you can also use these to access the fonts in for instance special
situations, like when you construct a title page: 

\starttyping
\definefont[TitleFont][Sans at 72pt]
\definefont[TitleFont][OfficinaSans-Book at 72pt]       
\stoptyping

Although you can also access the file name directly, these methods are more
descriptive. 


\subject {Tweaking shapes}

A font is seldom just one shape and file. Apart from the upright shape, there
can be an italic, bold and bold||italic alternative, or, oblique, bold and
bold||oblique. Italic and oblique share their forward slanted look. By
applying appropriate transformations, you can slant any font or make it wider
or narrower. 

So, we have at our hands, either or not by manipulation, normal, italic or
oblique, slanted, as well as their bold variants. The following command
creates a slanted bold officina font. 

\starttyping
texfont --ve=itc --co=officina --sl=* ovb 
\stoptyping

Instead of a \type {*}, you can provide a number. The default slant is
$.167$. Another manipulation is to widen a font, using the extend key: 

\starttyping
texfont --ve=itc --co=officina --ex=1.2 ovb 
\stoptyping

These commands lead to font metric files with names as: 

\starttyping
texnansi-ovb-slanted-167.tfm 
texnansi-ovb-extended-1200.tfm 
\stoptyping

Combinations are also possible. In the default \TEX\ distributions, where the
8 character file name limit is still honoured, a less verbose naming scheme
is used. Our alternative is less efficient, but opens the possibility to have
more then one slanted alternative. 

Since \TEXFONT\ is mainly a wrapper around \type {afm2tfm}, we also provide a
third manipulation: creating small caps fonts. 

\starttyping
texfont --ve=itc --co=officina --ca=* ovb 
\stoptyping

We now get: 

\starttyping
texnansi-ovb-capitalized-800.tfm 
\stoptyping

\subject {Generating instances}

The Computer Modern Roman typefaces were created by Donald Knuth and their
shapes are described as programs. There are quite some design axis and
parameters that can be tweaked in order to get different instances. This is
why they qualify as meta||fonts. In its own way, Adobe has created the
Multiple Master Fonts format. \footnote {Probably due to lack of interest,
Adobe is no longer advocating this format, although it will survive in Open
Type fonts.} 

If you decide (or are forced) to use a multiple master font, you need to make
sure that you get all the files needed to tweak them. For instance, for using
the Myriad fonts, you need files like: 

\starttyping
MyriadMM-LightCn.afm            MyriadMM-LightCnIt.afm
MyriadMM-BlackCn.afm            MyriadMM-BlackCnIt.afm
MyriadMM-LightSemiEx.afm        MyriadMM-LightSemiExIt.afm
MyriadMM-BlackSemiEx.afm        MyriadMM-BlackSemiExIt.afm
MyriadMM.amfm                   MyriadMM-It.amfm
MyriadMM.pfb                    MyriadMM-It.pfb
\stoptyping

When the author bought this font, it came as an install binary, that needed
the Adobe Type Manager. After a couple of failures, he finally found some
\type {pfb} and \type {mmm} files on my system. 

You need the programs \type {mmafm} and \type {mmpfb} to create the
instances. \footnote {Currently these are only available for the Unix
operating system.} These programs need an \type {amfm} file, which did not
came on the floppy, but a bit of emailing and renaming finally lead to the
files mentioned previously. 

The following commands will create an acceptable collection of normal and
italic, bold and bold italic shapes. 

\starttyping
texfont --we=400 --wi=600 MyriadMM
texfont --we=700 --wi=600 MyriadMM
texfont --we=400 --wi=600 MyriadMM-It
texfont --we=700 --wi=600 MyriadMM-It
\stoptyping

After this you will have files like: 

\starttyping
MyriadMM-weight-400-width-600.afm 
MyriadMM-weight-400-width-600.pfb 
\stoptyping

These alternatives can now be made \TEX||ready by saying: 

\starttyping
texfont --ve=adobe --co=myriad --ma --in MyriadMM-we*
texfont --ve=adobe --co=myriad --ma --in MyriadMM-It-we*
\stoptyping

Slanted, boldslanted and capitalized variants can be
created with: 

\starttyping
texfont --ve=adobe --co=myriad --sl=* MyriadMM-we*
texfont --ve=adobe --co=myriad --ca=* MyriadMM-we*
\stoptyping

\subject {Map files}

There are several ways to tell \PDFTEX\ which map files to 
use. One way is to add an entry to the file \type 
{pdftex.cfg}, like:

\starttyping
map +texnansi-urw-palatino.map 
\stoptyping

You can also add an entry in the \TEX\ file, by saying: 

\starttyping
\pdfmapfile{+texnansi-urw-palatino.map} 
\stoptyping

or in \CONTEXT\ with: 

\starttyping
\loadmapfile[texnansi-urw-palatino.map]
\stoptyping

This command prevents duplicate loading of files. If you want \CONTEXT\ to 
load the files automatically, you can add an entry to your \type 
{cont-sys.tex} file:

\starttyping
\autoloadmapfilestrue
\stoptyping

Beware: currently \PDFTEX\ only loads map files before the first page is 
shipped out. If you define fonts halfway the document, you must make sure 
that the associated map file is loaded at the top of your document. A rather 
rough way out is to say: 

\starttyping
\usetypescript[map][texnansi,ec,8r]
\stoptyping

or whatever combination of encodings you want. If you provide the keyword 
\type {all}, \CONTEXT\ will load all known map files. 

\subject {Switches}

You can control \TEXFONT's behaviour with command line switches. You can get
an overview of these by saying 

\starttyping
texfont --help
\stoptyping

Some switches expect a number or string. You can abbreviate switches as long
as you make sure that they can be distinguished. 

\switch fontroot=path

This is the place where the files that are generated will go to. You can
either set the font root manually, or let \TEXFONT\ consult your path
settings. 

\switch sourcepath=path 

When you install new fonts, by default the current path is taken. This switch
can be used to specify an absolute or relative path. If you provide \type
{auto} as path, \TEXFONT\ will try to locate the source path by means of the
vendor and collection specification. 

\switch vendor=name 

When generating metrics, this key is mandatory. It is used as a directory
specifier under the \type {/fonts} path and ends up as part of the map file
name. 

\switch collection=name 

Like the \type {vendor} key, when generating metrics, this key is mandatory.
It is used as a directory specifier under the \type {/vendor} path and ends
up as part of the map file name. 

\switch encoding=name 

Here you specify the font encoding vector. You need to make sure that you
have a correct file on your system (like \type {ec.enc} or \type {8r.enc}).
The default encoding is \type {texnansi} but any decent alternative will do. 

\switch slant=number 

The number specified here is normally not that large. An often used value is
\type {.167} which is also the default. If you provide a \type {*}, the
default value is used. Values between $0.0$ and $1.5$ are accepted. The
number, multiplied by $1000$ and rounded ends up in the name of the font
instance. 

\switch extend=number     

This specifier determines how much a glyph will be stretched in horizontal
direction. The default value is $1.2$. See also \type {slant}. 

\switch caps=number 

This specifier determines how much lowercase glyphs will be scaled in
vertical direction. The default value is $0.8$. See also \type {slant}. 

\switch weight=number 

When generating a multiple master instance, this number determines how bold a
glyph will look. Values between 200 and 700 give acceptable results. The
exact specification depends on the font. The number becomes part of the
filename. 

\switch width=number 

This number defines how wide a glyph will be. See also \type {weight}. 

\switch install           

This switch instructs \TEXFONT\ to copy files from the source path into the
right locations of the font tree. 

\switch makepath          

If you provide this switch, \TEXFONT\ will create the paths needed for
installing the fonts. Otherwise, you have to do so yourself, otherwise
\TEXFONT\ will quit with a warning. 

\switch listing 

Given that your sourcepath is set to \type {auto}, this switch will result in
a list of the metric files of the installed fonts. 

\switch remove 

Given that your sourcepath is set to \type {auto}, this switch will result in
deletion of the metric files of the installed fonts. 

\switch test 

If you are just playing a bit with \TEXFONT, you may not want to clobber your
font path with files that you will never use. This switch sets the vendor and
collection equal to \type {test}. 

\switch show 

If you want, \TEXFONT\ can process the test file it generates during the
installation. 

\switch batch   

In the next section we will discuss batch files. This switch instructs
\TEXFONT\ to treat the filename given as a batch file. 

If you peek into the source of \TEXFONT, you will notice a couple of more
switches. These are not documented here, and may disappear in future
versions. 

\switch virtual 

Create a virtual font (\type {vf} and \type {tfm} file) instead of a normal 
one (\type {tfm} only). 

% passon=string
% virtual 
% narrow 
% noligs 

\subject {Batch files}

Once a decent conversion is sorted out, keying in commands like the ones
mentioned in previous sections will become an annoyance. This is why
\TEXFONT\ supports a crude but effective way of processing batch files. One
of the batch files that comes with \TEXFONT, \type {type-tmf.dat}, has
entries like: 

\starttyping
--en=? --ve=urw --co=palatino --so=auto
--en=? --ve=urw --co=palatino --so=auto --ca=* uplr8a
--en=? --ve=urw --co=palatino --so=auto --sl=* uplr8a
--en=? --ve=urw --co=palatino --so=auto --sl=* uplb8a
\stoptyping

The \type {?} will be replaced by the encoding passed to \TEXFONT\ when
processing this batch file: 

\starttyping
texfont --encoding=8r --batch type-tmf.dat 
\stoptyping

If you have \TEX\ Live installed, you can play with this file and the
examples shown in this batch file, since it only uses the free fonts that are
on the system. If you run the whole file, for which you can best set up a
font tree first, you get a nice collection of fonts. This way you can build
your own (stable) font tree. Another batch file is \type {type-buy.dat}, a 
file used at \PRAGMA. This file has a companion \type {type-buy.tex} file 
that contains typescripts. You need to load this file explicitly in order to 
get access to these scripts.

\starttyping
\usetypescriptfile[type-buy] 
\stoptyping

If you run a batch file with \type {vendor} and \type {collection} keys, only
the section marked as such will be processed. A section is marked with:
\footnote {Instead of \type {#}, a \type {%} may be used.} 

\starttyping
# urw palatino 
\stoptyping 

If you make batch files (or sections) for each font collection that you
install, you can always regenerate the files quickly. That way you can
occasionally clean up your disk and make a fresh start. 

On the next pages we will show the public Zapf Chancery font in some rather 
common encodings. Watch out: the current release of \PDFTEX\ only reads map 
files that are defined before the first page of the document is shipped out. 

\page 

\starttyping 
texfont --encoding=texnansi --vendor=urw --collection=zapfchan
\stoptyping 

\startbuffer
\loadmapfile[texnansi-urw-zapfchan]

\definefontsynonym[ZapfChancery][texnansi-uzcmi8a][encoding=texnansi] 

\showfont[ZapfChancery] \showligatures[ZapfChancery]
\stopbuffer

\typebuffer \getbuffer \page 

\starttyping 
texfont --encoding=8r --vendor=urw --collection=zapfchan
\stoptyping 

\startbuffer
\loadmapfile[8r-urw-zapfchan]

\definefontsynonym[ZapfChancery][8r-uzcmi8a][encoding=8r] 

\showfont[ZapfChancery] \showligatures[ZapfChancery]
\stopbuffer

\typebuffer \getbuffer \page 

\starttyping 
texfont --encoding=ec --vendor=urw --collection=zapfchan
\stoptyping 

\startbuffer
\loadmapfile[ec-urw-zapfchan]

\definefontsynonym[ZapfChancery][ec-uzcmi8a][encoding=ec]

\showfont[ZapfChancery] \showligatures[ZapfChancery]
\stopbuffer

\typebuffer \getbuffer \page

\ColofonPage

\stoptext 
